\documentclass[draft,jgrga]{agutexSI2019}

\usepackage{url}
\usepackage{amsmath}

\begin{document}

\title{Supporting Information for ``Aerosols in the vacuum: modelling the time variability of the plume of Enceladus''}

\authors{=list all authors here=}
\affiliation{=number=}{=Affiliation Address=}

\correspondingauthor{=name=}{=email address=}

\begin{article}

\noindent\textbf{Contents of this file}
\begin{enumerate}
\item Text S1 to S3
\end{enumerate}

\noindent\textbf{Introduction}

This Supporting Information collects the detailed derivations referenced in the
main text: the depth-integrated momentum equation that governs the liquid column
(Text S1), the energy balance of the thin wall thermal layer that sets the
wall sublimation/condensation rate (Text S2), and the quasi-one-dimensional
relation that governs the Mach number of the gas column (Text S3).

%%%%%%%%%%%%%%%%%%%%%%%%%%%%%%%%%%%%%%%%%%%%%%%%
\noindent\textbf{Text S1. Depth-integrated momentum equation for the liquid column}

We use a column coordinate $z$ running from the crack floor ($z=0$) to the
gas--liquid interface ($z=h+L$), where $L$ is the equilibrium water-column depth
and $h$ the displacement of the interface from equilibrium. The crack half-gap
$\delta(t)$ is prescribed (tidal forcing); a dot denotes $d/dt$.

\paragraph{Velocity profile from continuity.}
For an incompressible column in a slot of uniform gap $\delta(t)$, volume
conservation per unit crack length, $\partial\delta/\partial t +
\partial(\delta v)/\partial z = 0$, integrates from the floor (where $v=v_0$) to
\begin{equation}
    v(z,t) = v_0(t) - \frac{\dot\delta}{\delta}\,z,
    \qquad 0 \le z \le h+L,
    \label{eq:si-vel}
\end{equation}
which is Equation~(2) of the main text. The velocity at the interface is
$v_h \equiv v(h+L,t) = v_0 - (\dot\delta/\delta)(h+L)$, giving the kinematic
condition $dh/dt = v_h$ (Equation~(6) of the main text).

\paragraph{Pointwise momentum equation.}
Dividing the momentum balance (Equation~(3) of the main text) by $\rho_w\delta$,
\begin{equation}
    \frac{\partial v}{\partial t}
    = -\,v\frac{\partial v}{\partial z}
      - \frac{1}{\rho_w}\frac{\partial p}{\partial z}
      - g
      - \frac{2 C_f}{\delta}\,v^2\hat v,
    \label{eq:si-mom}
\end{equation}
where the last term is the friction of the two walls (stress $C_f\rho_w v^2$ on
each, hydraulic factor $2/\delta$) and $\hat v=\mathrm{sgn}(v)$.

\paragraph{Depth integration.}
Integrate \eqref{eq:si-mom} from $z=0$ to $z=h+L$. The advective term gives
\begin{equation}
    \int_0^{h+L} v\,\frac{\partial v}{\partial z}\,dz
    = \tfrac12\!\left(v_h^2 - v_0^2\right).
\end{equation}
The pressure term uses the two boundary conditions: a fixed hydrostatic pressure
$p(0)=\rho_w g L$ at the floor and $p(h+L)=0$ at the free surface, so
\begin{equation}
    -\frac{1}{\rho_w}\int_0^{h+L}\frac{\partial p}{\partial z}\,dz
    = -\frac{1}{\rho_w}\big[p(h+L)-p(0)\big] = gL,
\end{equation}
which combines with the gravity term $-g(h+L)$ to give $-gh$. Writing the
friction integral as $F_{\rm fric}=\int_0^{h+L}(2C_f/\delta)v^2\hat v\,dz$, the
right-hand side of the integrated equation is
\begin{equation}
    \int_0^{h+L}\frac{\partial v}{\partial t}\,dz
    = -\tfrac12\!\left(v_h^2-v_0^2\right) - gh - F_{\rm fric}.
    \label{eq:si-rhs}
\end{equation}

\paragraph{Left-hand side.}
From \eqref{eq:si-vel}, the explicit time dependence of the profile is
\begin{equation}
    \frac{\partial v}{\partial t}
    = \frac{dv_0}{dt}
      - \frac{d}{dt}\!\left(\frac{\dot\delta}{\delta}\right) z
    = \frac{dv_0}{dt}
      + \left[\left(\frac{\dot\delta}{\delta}\right)^{\!2} - \frac{\ddot\delta}{\delta}\right] z ,
\end{equation}
using $\tfrac{d}{dt}(\dot\delta/\delta)=\ddot\delta/\delta-(\dot\delta/\delta)^2$.
Integrating over the column,
\begin{equation}
    \int_0^{h+L}\frac{\partial v}{\partial t}\,dz
    = (h+L)\,\frac{dv_0}{dt}
      + \frac12\left[\left(\frac{\dot\delta}{\delta}\right)^{\!2} - \frac{\ddot\delta}{\delta}\right](h+L)^2 .
    \label{eq:si-lhs}
\end{equation}

\paragraph{Result.}
Equating \eqref{eq:si-lhs} and \eqref{eq:si-rhs} and solving for $dv_0/dt$,
\begin{equation}
    \frac{dv_0}{dt}
    = \frac{1}{h+L}\left[
      -\tfrac12\!\left(v_h^2-v_0^2\right)
      - \tfrac12\!\left[\left(\frac{\dot\delta}{\delta}\right)^{\!2} - \frac{\ddot\delta}{\delta}\right](h+L)^2
      - gh - \int_0^{h+L}\frac{2C_f}{\delta}v^2\hat v\,dz
    \right],
    \label{eq:si-dv0}
\end{equation}
which is Equation~(4) of the main text. The factor $(h+L)^2$ on the second
(unsteady) term follows from the depth integration in \eqref{eq:si-lhs} and is
required dimensionally; the advection term carries the sign $-\tfrac12(v_h^2-v_0^2)$.

For a height-dependent gap $\delta(z,t)$ (a tapered crack), the same procedure
applies with $v(z)=[\delta(z_0)v_0-\int_{z_0}^z\partial_t\delta\,dz']/\delta(z)$;
the inertial pre-factor $(h+L)$ in \eqref{eq:si-dv0} is then replaced by
$\int_0^{h+L}\delta(z_0)/\delta(z)\,dz$, and the unsteady term by the
corresponding integral of the explicit profile time-dependence.

%%%%%%%%%%%%%%%%%%%%%%%%%%%%%%%%%%%%%%%%%%%%%%%%
\noindent\textbf{Text S2. Energy balance of the wall thermal layer}

The thin wall layer (region 5, thickness $\lambda\simeq1$~m, several diurnal
skin depths) connects the in-crack heat exchange to the conductive interior.
Its surface temperature $T_s$ at a point a distance $d$ below the crack mouth is
set by balancing the radiative exchange with the cold sky and the conductive
supply from the interior. Treating the surface as grey with the sky at the
equilibrium temperature $T_e$ and the conductive flux from the interior
(at base temperature $T_b$) over a path of order $d$ (Equation~(8) of the main
text, $F_c = 4k\,[T_s-T_b]/(\pi d)$), the steady balance is
\begin{equation}
    2\sigma\!\left(T_s^4 - T_e^4\right) + \frac{4k}{\pi d}\left(T_s - T_b\right) = 0 ,
    \label{eq:si-Ts}
\end{equation}
solved for the positive root $T_s(d)$ (a bisection in the code). The net
sublimation rate per unit wall area follows from the surface energy released,
\begin{equation}
    E(d) = -\frac{\sigma}{L_v}\left(T_s^4 - T_e^4\right),
    \label{eq:si-E}
\end{equation}
with $E<0$ denoting condensation. Two limits are instructive: near the mouth
($d\to0$) the conduction term dominates, $T_s\to T_b$, and $|E|$ is largest
(condensation concentrates within $\sim$one skin depth of the surface); deep in
the crack ($d$ large) the conduction term vanishes, $T_s\to T_e$, and $E\to0$.

When a wall element is instead submerged, it exchanges heat with the liquid by
conduction across region~5; the time-dependent diffusion equation
$\partial_t T = \kappa\,\partial_{zz}T$ is integrated in the layer with the
sublimation-side boundary set by \eqref{eq:si-E} (when exposed) or by contact
with the $\sim$273~K liquid (when submerged), and the conduction-side boundary
matched to the steady interior profile of Equation~(8). The saturated vapour
pressure entering the exposed-wall flux uses
$p_w = A\exp(-B/T_w)$ with $A=3.63\times10^{12}$~Pa, $B=6147$~K.

%%%%%%%%%%%%%%%%%%%%%%%%%%%%%%%%%%%%%%%%%%%%%%%%
\noindent\textbf{Text S3. Mach-number relation for the gas column}

The vapour column (region~2) is modelled as quasi-one-dimensional, compressible,
calorically perfect flow in a slot of gap $\delta$, subject to distributed mass
addition from wall sublimation ($E$ per unit wall area) and wall friction
($C_f$). Combining the differential forms of continuity, momentum, energy and
the perfect-gas relation for such a flow (the influence-coefficient method),
the Mach number $M$ evolves with height $z$ as
\begin{equation}
    \frac{dM}{M}\,\frac{1-M^2}{1+\frac{\gamma-1}{2}M^2}
    = \left(\frac{E}{\rho v \delta} + \frac{\gamma C_f}{2\delta}M^2\right)dz ,
    \label{eq:si-mach}
\end{equation}
which is Equation~(9) of the main text. The first source term is mass addition:
$E/(\rho v\delta)$ is the fractional rate of mass injection per unit length from
the two walls into the column of cross-section $\propto\delta$. The second is
the Fanning wall friction, $(\gamma C_f/2\delta)M^2$, with the factor $1/\delta$
the hydraulic length of the slot. As in adiabatic duct flow, the
$\left(1-M^2\right)$ factor makes the two sources drive $M$ toward unity from
either side, so the exit Mach number is insensitive to the precise boundary
value over the range $1.4$--$1.8$ inferred for the plume; we adopt $M=1.6$ at the
crack mouth.

\end{article}
\end{document}
